\SourcePage{419}
\SGAchapter{XII}{代数几何中的 Nielsen--Wecken 公式与 Lefschetz 公式}{Formules de Nielsen--Wecken et de Lefschetz en géométrie algébrique}
\begin{center}A. Grothendieck\quad 著\\I. Bucur\quad 整理\end{center}

\SourcePage{420}
\SGAsection{1}{基本设定}
设 $X'$ 是\Term{algebraically-closed-field}{代数闭域} $k$ 上的\Term{proper}{固有}\Term{scheme}{概形}，$G$ 是\Term{action}{作用}于 $X'$ 的\Term{finite-group}{有限群}。对于每个 $g\in G$，以 $g_{X'}:X'\longrightarrow X'$ 表示 $g$ 所诱导的 $X'$ 的\Term{automorphism}{自同构}。

假定 $X'$ 有一个在 $G$ 下\Term{stable}{稳定}的\Term{open-subscheme}{开子概形} $U'$，使得 $G$ 在 $U'$ 上\Term{free-action}{自由作用}。令 $\Lambda_n=\mathbf Z/\ell^n\mathbf Z$，其中 $\ell$ 是异于 $k$ 的\Term{characteristic}{特征}的素数。

以 $\Lambda_{n,U',X'}$ 表示 $X'$ 上的\Term{sheaf}{层} $i_!(\Lambda_{n,U'})$，其中 $\Lambda_{n,U'}$ 是 $U'$ 上与环 $\Lambda_n$ 对应的\Term{constant-sheaf}{常值层}，$i:U'\hookrightarrow X'$ 是\Term{canonical}{典范}\Term{immersion}{浸入}。\Term{sheaf}{层} $\Lambda_{n,U',X'}$ 自然带有以 $G$ 为\Term{operator-group}{算子群}的\Term{sheaf}{层}结构。因此，\Term{complex}{复形}
\[
\mathbb H_{n,U',X'}=\mathrm R\Gamma_{X'}(\Lambda_{n,U',X'})
\]
是 $\Lambda_n[G]$-\Term{module}{模}\Term{complex}{复形}；更确切地说，它是 $\Lambda_n[G]$-\Term{module}{模}的\Term{perfect-complex}{完美复形}（第 VIII 讲 7；第 IX 讲 2）。

我们将看到，为达到这里的目的，使用 $\mathrm R\Gamma_{X'}(\Lambda_{n,U',X'})$ 相对于环 $\Lambda_n$ 在\Term{derived-category}{导出范畴}意义下的\Term{dual}{对偶}，比使用它本身更方便。记此\Term{dual}{对偶}为
\[
\begin{aligned}
\mathbb H_{n,U',X'}^\vee
 &=\mathrm R\Gamma_{X'}(\Lambda_{n,U',X'})^\vee\\
 &=\mathrm R\operatorname{Hom}_{\Lambda_n}
   \bigl(\mathrm R\Gamma_{X'}(\Lambda_{n,U',X'}),\Lambda_n\bigr).
\end{aligned}
\]
\Term{complex}{复形} $\mathrm R\Gamma_{X'}(\Lambda_{n,U',X'})^\vee$ 也是 $\Lambda_n[G]$-\Term{module}{模}的\Term{perfect-complex}{完美复形}。

现在再给定\Term{scheme}{概形} $X'$ 的一个\Term{endomorphism}{自同态} $f':X'\longrightarrow X'$，以及群 $G$ 的一个\Term{endomorphism}{自同态} $\varphi:G\longrightarrow G$，使以下两个条件成立：
\SourcePage{421}
\[
f'g_{X'}=\varphi(g)_{X'}f'\qquad(\forall g\in G).
\tag{1.1}
\]
\[
\text{补集 }Y'=X'-U'\text{ 在 }f'\text{ 下稳定。}
\tag{1.2}
\]
（由已有假设，$Y'$ 在 $G$ 下\Term{stable}{稳定}。一般并不要求\Term{open-subscheme}{开子概形} $U'$ 在 $f'$ 下\Term{stable}{稳定}。）

于是得到下列\Term{sheaf}{层}\Term{homomorphism}{同态}：
\[
\begin{gathered}
f'_{\Lambda_{n,X'}}:\Lambda_{n,X'}\longrightarrow\Lambda_{n,X'},\qquad
f'_{\Lambda_{n,Y'}}:\Lambda_{n,Y'}\longrightarrow\Lambda_{n,Y'},\\
f'_{\Lambda_{n,U',X'}}:\Lambda_{n,U',X'}\longrightarrow\Lambda_{n,U',X'},\qquad
g_{\Lambda_{n,U',X'}}:\Lambda_{n,U',X'}\longrightarrow\Lambda_{n,U',X'},
\end{gathered}
\]
并且下图交换，其中各行均\Term{exact}{正合}：
\[
\begin{tikzcd}[column sep=small]
0\arrow[r]&\Lambda_{n,U',X'}\arrow[r]\arrow[d,"{f'_{\Lambda_{n,U',X'}}}"']&\Lambda_{n,X'}\arrow[r]\arrow[d,"{f'_{\Lambda_{n,X'}}}"]&\Lambda_{n,Y'}\arrow[r]\arrow[d,"{f'_{\Lambda_{n,Y'}}}"]&0\\
0\arrow[r]&\Lambda_{n,U',X'}\arrow[r]&\Lambda_{n,X'}\arrow[r]&\Lambda_{n,Y'}\arrow[r]&0.
\end{tikzcd}
\tag{1.3}
\]
特别地，由此得到下列 $\Lambda_n$-\Term{module}{模}\Term{complex}{复形}的\Term{homomorphism}{同态}：
\SourcePage{422}
\[
\begin{aligned}
f'_{\mathbb H_{n,U',X'}}
 &=\mathrm R\Gamma_{X'}(f'_{\Lambda_{n,U',X'}}):
 \mathbb H_{n,U',X'}\longrightarrow\mathbb H_{n,U',X'},\\
f'_{\mathbb H_{n,U',X'}^\vee}
 &=(f'_{\mathbb H_{n,U',X'}})^\vee:
 \mathbb H_{n,U',X'}^\vee\longrightarrow\mathbb H_{n,U',X'}^\vee,\\
g_{\mathbb H_{n,U',X'}}
 &=\mathrm R\Gamma_{X'}(g_{\Lambda_{n,U',X'}}):
 \mathbb H_{n,U',X'}\longrightarrow\mathbb H_{n,U',X'},\\
g_{\mathbb H_{n,U',X'}^\vee}
 &=(g_{\mathbb H_{n,U',X'}})^\vee:
 \mathbb H_{n,U',X'}^\vee\longrightarrow\mathbb H_{n,U',X'}^\vee.
\end{aligned}
\]
它们满足下列关系：
\[
g_{\mathbb H_{n,U',X'}}\circ f'_{\mathbb H_{n,U',X'}}
=f'_{\mathbb H_{n,U',X'}}\circ\varphi(g)_{\mathbb H_{n,U',X'}};
\tag{1.4}
\]
\[
f'_{\mathbb H_{n,U',X'}^\vee}\circ g_{\mathbb H_{n,U',X'}^\vee}
=\varphi(g)_{\mathbb H_{n,U',X'}^\vee}\circ f'_{\mathbb H_{n,U',X'}^\vee}.
\tag{1.5}
\]
由 (1.5) 可知（参见关于迹的一般理论的第 XI 讲 5.2.1、7.1），其\Term{trace}{迹}满足
\[
\operatorname{Tr}^{*}_{\Lambda_n[G],\varphi}
 (f'_{\mathbb H_{n,U',X'}^\vee})
\in\Lambda_n[G_{\natural,\varphi}]
\simeq\mathscr F_!(G_{\natural,\varphi},\Lambda_n).
\]
\footnote{参见第 III B 讲 5.12。}
我们的第一个目标，是在一些附加假设下，用与 $f'$ 相关的\Term{local-invariant}{局部不变量}给出
$\operatorname{Tr}^{*}_{\Lambda_n[G],\varphi}(f'_{\mathbb H_{n,U',X'}^\vee})$ 的显式公式；这些\Term{local-term}{局部项}将在下文定义。

\SGAsection{2}{对系数的相容性}
实际上还可以做得更多。为此，设 $n,m$ 是满足 $m\geq n$ 的两个正整数。环\Term{homomorphism}{同态}
\SourcePage{423}
\[
\alpha_{m,n}:\Lambda_m\longrightarrow\Lambda_n
\]
诱导 $X'$ 上的环\Term{sheaf}{层}\Term{homomorphism}{同态}
\[
\alpha_{m,n,X'}:\Lambda_{m,X'}\longrightarrow\Lambda_{n,X'},\qquad
\alpha_{m,n,U',X'}:\Lambda_{m,U',X'}\longrightarrow\Lambda_{n,U',X'}.
\]
应用 \Term{kunneth-formula}{Künneth 公式}，得到
\[
\mathrm R\Gamma_{X'}(\Lambda_{m,X'})\mathbin{\overset{\mathrm L}{\otimes}}_{\Lambda_m}\Lambda_n
\simeq\mathrm R\Gamma_{X'}(\Lambda_{n,X'}),
\tag{2.1}
\]
\[
\mathrm R\Gamma_{X'}(\Lambda_{m,U',X'})\mathbin{\overset{\mathrm L}{\otimes}}_{\Lambda_m}\Lambda_n
\simeq\mathrm R\Gamma_{X'}(\Lambda_{n,U',X'}),
\tag{2.2}
\]
\[
\mathrm R\Gamma_{X'}(f'_{\Lambda_{m,X'}})\mathbin{\overset{\mathrm L}{\otimes}}_{\Lambda_m}\Lambda_n
\simeq\mathrm R\Gamma_{X'}(f'_{\Lambda_{n,X'}}),
\tag{2.3}
\]
\[
\mathrm R\Gamma_{X'}(f'_{\Lambda_{m,U',X'}})\mathbin{\overset{\mathrm L}{\otimes}}_{\Lambda_m}\Lambda_n
\simeq\mathrm R\Gamma_{X'}(f'_{\Lambda_{n,U',X'}}).
\tag{2.4}
\]
容易由此\Term{pushout}{推出}
\[
\begin{aligned}
{}&\mathrm R\operatorname{Hom}_{\Lambda_m}
 (\mathrm R\Gamma_{X'}(\Lambda_{m,U',X'}),\Lambda_m)
 \mathbin{\overset{\mathrm L}{\otimes}}_{\Lambda_m}\Lambda_n\\
&\hspace{2em}\simeq\mathrm R\operatorname{Hom}_{\Lambda_n}
 (\mathrm R\Gamma_{X'}(\Lambda_{n,U',X'}),\Lambda_n),
\end{aligned}
\tag{2.5}
\]
\[
f'_{\mathbb H_{m,U',X'}^\vee}\mathbin{\overset{\mathrm L}{\otimes}}_{\Lambda_m}\Lambda_n
\simeq f'_{\mathbb H_{n,U',X'}^\vee}.
\tag{2.6}
\]
因此，\Term{perfect-complex}{完美复形}的\Term{morphism}{态射}
$\mathrm R\Gamma_{X'}(f'_{\Lambda_{n,X'}})$（以及相应的
$\mathrm R\Gamma_{X'}(f'_{\Lambda_{n,U',X'}})$、$f'_{\mathbb H_{n,U',X'}^\vee}$）
由\Term{morphism}{态射} $\mathrm R\Gamma_{X'}(f'_{\Lambda_{m,X'}})$（以及相应的
$\mathrm R\Gamma_{X'}(f'_{\Lambda_{m,U',X'}})$、$f'_{\mathbb H_{m,U',X'}^\vee}$）
沿 $\alpha_{m,n}$（相应地沿 $\alpha_{m,n}$、$\alpha_{m,n,G}:\Lambda_m[G]\longrightarrow\Lambda_n[G]$）作\Term{extension-of-scalars}{标量扩张}得到（第 XI 讲 6.3）。
\SourcePage{424}
利用第 XI 讲 6.3、6.3.2、6.3.3 以及 (2.3)、(2.6)，得到
\[
\operatorname{Tr}^{*}_{\Lambda_n}(\mathrm R\Gamma_{X'}(f'_{\Lambda_{n,X'}}))
=\alpha_{m,n}\bigl(\operatorname{Tr}^{*}_{\Lambda_m}
 (\mathrm R\Gamma_{X'}(f'_{\Lambda_{m,X'}}))\bigr),
\tag{2.7}
\]
\[
\operatorname{Tr}^{*}_{\Lambda_n[G],\varphi}(f'_{\mathbb H_{n,U',X'}^\vee})
=\alpha_{m,n,G}\bigl(\operatorname{Tr}^{*}_{\Lambda_m[G],\varphi}
 (f'_{\mathbb H_{m,U',X'}^\vee})\bigr).
\tag{2.8}
\]
函数族
$\bigl(\operatorname{Tr}_{\Lambda_n[G],\varphi}(f'_{\mathbb H_{n,U',X'}^\vee})\bigr)_{n>0}$
确定了 $G$ 上的一个取值于\Term{ell-adic-integer}{$\ell$ 进整数}环 $\mathbf Z_\ell$ 的
\Term{twisted-central-function}{$\varphi$ 中心函数}
\[
\operatorname{Tr}_{\mathbf Z_\ell[G],\varphi}
 (f'_{\mathbb H_{U',X'}^\vee(\ell)})
\in\mathscr F_!(G_{\natural,\varphi},\mathbf Z_\ell).
\]
同样，元素
$\operatorname{Tr}^{*}_{\Lambda_n}(\mathrm R\Gamma_{X'}(f'_{\Lambda_{n,X'}}))$，$n\geq1$，
确定了一个\Term{ell-adic-integer}{$\ell$ 进整数}
\[
\operatorname{Tr}^{*}_{\mathbf Z_\ell}
 (f'_{\mathrm R\Gamma_{X'}(\mathbf Z_{\ell,X'})}).
\]
我们的目标就是求出
$\operatorname{Tr}^{*}_{\mathbf Z_\ell[G],\varphi}(f'_{\mathbb H_{U',X'}^\vee(\ell)})$ 的公式。

\SGAsection{3}{Nielsen--Wecken 局部不变量}
参见文献 [1]。在第 1 节的假设之外，再假定 $G$ 的每条\Term{orbit}{轨道}都包含于某个\Term{affine}{仿射}\Term{open-subscheme}{开子概形}中。于是\Term{quotient-scheme}{商概形} $X=X'/G$ 存在；根据 (1.1)，$X'$ 的\Term{endomorphism}{自同态} $f'$ 可\Term{descent}{下降}为\Term{quotient-scheme}{商概形}上的\Term{endomorphism}{自同态}
\[
f:X\longrightarrow X.
\]
下文使用以下记号：
\SourcePage{425}
\[
p:X'\longrightarrow X=X'/G\quad\text{为典范投影},\qquad
Y=p(Y')=Y'/G,\qquad U=X-Y.
\]
$U'$ 是以 $G$ 为\Term{structure-group}{结构群}的\Term{principal-covering}{主覆盖}。设 $x\in U$ 是 $f$ 的一个\Term{fixed-point}{不动点}。我们按以下方式给它关联一个
\Term{twisted-conjugacy-class}{$\varphi$ 共轭类}（第 XI 讲 7.1）$\operatorname{NW}_{f'}^{G,\varphi}(x)$。

取 $x'\in U'$ 使 $p(x')=x$。由于 $U'$ 是以 $G$ 为\Term{structure-group}{结构群}的\Term{principal-covering}{主覆盖}，且 $x$ 在 $f$ 下不动，存在唯一的 $g\in G$，使得 $f'(x')=gx'$。$g$ 的\Term{twisted-conjugacy-class}{$\varphi$ 共轭类}只依赖于 $x$。事实上，若 $p(x'')=x$，则存在 $h\in G$ 使 $hx''=x'$。因此，若 $f'(x'')=g'x''$，便有
\[
gx'=f'(x')=f'(hx'')=\varphi(h)f'(x'')=(\varphi(h)g'h^{-1})x',
\]
\[
g=\varphi(h)g'h^{-1}.
\]
也就是说，$g$ 与 $g'$ 是 $\varphi$ \Term{conjugacy}{共轭}的。由此得到一个确定的\Term{twisted-conjugacy-class}{$\varphi$ 共轭类}，称为与 $f$ 的\Term{fixed-point}{不动点} $x\in U$ 相关联的
\Term{nielsen-wecken-invariant}{Nielsen--Wecken 局部不变量}，记作 $\operatorname{NW}_{f'}^{G,\varphi}(x)$。

\Term{twisted-conjugacy-class}{$\varphi$ 共轭类} $\operatorname{NW}_{f'}^{G,\varphi}(x)$ 确定一个整数值的
\Term{twisted-central-function}{$\varphi$ 中心函数}（参见第 XI 讲 7.1）
\[
\operatorname{NW}_{f'}^{G,\varphi}(x):G\longrightarrow\mathbf Z,
\]
它由以下关系刻画：
\SourcePage{426}
\[
\operatorname{NW}_{f'}^{G,\varphi}(x)(g)=
\begin{cases}
1,&\text{存在 }x'\in X'\text{，使 }(g^{-1}f')(x')=x'\text{ 且 }p(x')=x,\\
0,&\text{否则。}
\end{cases}
\tag{3.1}
\]

\SGAsection{4}{局部函数及其整性}
假定 $X'$ 在 $k$ 上\Term{smooth}{光滑}，并且 $f$ 的\Term{fixed-point}{不动点}个数\Term{finite}{有限}。特别地，对每个 $g\in G$，$g^{-1}f'$ 的\Term{fixed-point}{不动点}个数也\Term{finite}{有限}。

对于 $f$ 的每个\Term{fixed-point}{不动点} $x\in X$，引入一个取有理数值的
\Term{twisted-central-function}{$\varphi$ 中心函数}
\[
S_x(f',f,G,\varphi):G\longrightarrow\mathbf Q,
\]
定义为
\[
S_x(f',f,G,\varphi)(g)
=\frac{1}{c^\varphi(g)}\sum_{x'\in X_x'{}^{g^{-1}f'}}
 \bigl(\nu_{x'}(g^{-1}f')-1\bigr).
\tag{4.1}
\]
这里使用了以下记号：
\[
c^\varphi(g)=\operatorname{Card}G_g^\varphi,\qquad
G_g^\varphi=\{a\in G\mid\varphi(a)ga^{-1}=g\},
\tag{4.2}
\]
\[
X_x'{}^{g^{-1}f'}
=\{x'\in X'\mid p(x')=x,\ (g^{-1}f')(x')=x'\},
\tag{4.3}
\]
\[
\nu_{x'}(g^{-1}f')=g^{-1}f'\text{ 的不动点 }x'\text{ 的重数。}
\tag{4.4}
\]
将上述函数与\Term{canonical}{典范}\Term{embedding}{嵌入} $\mathbf Q\hookrightarrow\mathbf Q_\ell$ 复合，
\SourcePage{427}
便由 $S_x(f',f,G,\varphi)$ 得到 $G$ 上取值于\Term{ell-adic-number}{$\ell$ 进数}域 $\mathbf Q_\ell$ 的
\Term{twisted-central-function}{$\varphi$ 中心函数} $S_{x,\ell}(f',f,G,\varphi)$。

\begin{SGAresult}{猜想}{4.5}
对每个 $\ell\ne\operatorname{char}(k)$，函数 $S_{x,\ell}(f',f,G,\varphi)$ 的值都是\Term{ell-adic-integer}{$\ell$ 进整数}。
\end{SGAresult}
下面将证明，当 $X$ 的\Term{dimension}{维数}等于 $1$ 时，猜想 4.5 成立。

\begin{SGAresult}{命题}{4.6}
设 $R$ 是\Term{discrete-valuation-ring}{离散赋值环}，$\mathfrak m$ 是它的\Term{maximal-ideal}{极大理想}，$\pi$ 是一个\Term{uniformizer}{素元}，$k=R/\mathfrak m$ 是它的\Term{residue-field}{剩余域}，并假定 $k$ 为\Term{algebraically-closed-field}{代数闭域}。设 $v$ 是它的\Term{valuation}{赋值}，$f'_R:R\longrightarrow R$ 是 $R$ 的\Term{endomorphism}{自同态}，$G$ 是 $R$ 的一个\Term{finite}{有限}\Term{automorphism}{自同构}群。假定 $f'_R$ 及各\Term{automorphism}{自同构} $g_R$ 在 $k$ 上都诱导\Term{identity-map}{恒等映射}，并且给定群\Term{endomorphism}{自同态}
\[
\varphi:G\longrightarrow G,
\]
使得
\[
f'_R\circ g_R=\varphi(g)_R\circ f'_R.
\]
对每个 $g\in G$，记
\[
\nu(g_R^{-1}f'_R)=v\bigl(g_R^{-1}f'_R(\pi)-\pi\bigr).
\]
将 $c^\varphi(g)$（参见 (4.2)）写为
\SourcePage{428}
\[
c^\varphi(g)=p^s q,\qquad(p,q)=1,\qquad p=\operatorname{char}(k),
\]
则 $q\mid\bigl(\nu(g_R^{-1}f'_R)-1\bigr)$。
\end{SGAresult}

为证明命题 4.6，可以归约到 $g=e$ 且 $\varphi=\operatorname{id}_G$ 的情形。事实上，先把情形 $(f'_R,G,g,\varphi)$ 换为
\[
(f''_R=g_R^{-1}f'_R,\ G,e,\varphi'=\operatorname{int}(g)\circ\varphi),
\qquad\varphi'(a)=g^{-1}\varphi(a)g.
\]
容易看到
\[
\nu(g_R^{-1}f'_R)-1=\nu(f''_R)-1,
\]
\[
c^\varphi(g)=c^{\varphi'}(e),\qquad G_g^\varphi=G_e^{\varphi'}
\quad\text{（参见 (4.2)）}.
\]
这就归约到 $g=e$。为了同时归约到 $g=e$、$\varphi=\operatorname{id}_G$，注意可以将
$(f''_R,G,e,\varphi')$ 换为
\[
(f''_R,G_e^{\varphi'},e,\operatorname{id}_{G_e^{\varphi'}}),
\]
并且有
\[
G_e^{\varphi'}=(G_e^{\varphi'})_e^{\operatorname{id}_{G_e^{\varphi'}}}.
\]
命题 4.6 的证明主要使用以下\Term{lemma}{引理}。

\begin{SGAresult}{\Term{lemma}{引理}}{4.7}
沿用命题 4.6 的假设与记号，再假定 $\varphi=\operatorname{id}_G$，即 $G$ 在 $R$ 上的\Term{action}{作用}与 $f'_R$ 交换，并且 $r=\nu(f'_R)\geq2$。令
\[
c=\operatorname{Card}\bigl(\operatorname{Im}
 (G\longrightarrow\operatorname{Aut}_k(\mathfrak m/\mathfrak m^2))\bigr).
\]
则有以下整除关系：
\SourcePage{429}
\[
c\mid(r-1).
\]
\end{SGAresult}
\begin{proof}
显然，整数 $r$ 满足下列条件：

\textup{a)} 由 $f'_R$ 诱导的\Term{homomorphism}{同态}
\[
f'_r:A/\mathfrak m^r\longrightarrow A/\mathfrak m^r
\]
是\Term{identity-map}{恒等映射}；

\textup{b)} \Term{homomorphism}{同态}
\[
f'_{r+1}:A/\mathfrak m^{r+1}\longrightarrow A/\mathfrak m^{r+1}
\]
不等于\Term{identity-map}{恒等映射}。

另一方面，$G$ 显然\Term{action}{作用}于所有的 $k$-\Term{vector-space}{向量空间} $\mathfrak m^i/\mathfrak m^{i+1}$。

我们首先证明：在\Term{lemma}{引理}的假设下，存在唯一一个 $k[G]$-\Term{module}{模}\Term{homomorphism}{同态}
\[
u(f'_R):\mathfrak m/\mathfrak m^2\longrightarrow
 \mathfrak m^r/\mathfrak m^{r+1},
\]
使以下两个条件成立：
\[
u(f'_R)\ne0;
\]
\[
u(f'_R)(x\bmod\mathfrak m^2)
=(f'_R(x)-x)\bmod\mathfrak m^{r+1}\qquad(\forall x\in\mathfrak m).
\]
显然，由条件 a)，对每个 $x\in R$ 都有 $f'_R(x)-x\in\mathfrak m^r$，因此第二个条件是有意义的。

\SourcePage{430}
这一观察还使我们能够定义一个\Term{abelian-group}{Abel 群}\Term{homomorphism}{同态}
\[
\mathfrak m\longrightarrow\mathfrak m^r/\mathfrak m^{r+1},
\qquad x\longmapsto(f'_R(x)-x)\bmod\mathfrak m^{r+1}.
\]
此\Term{homomorphism}{同态}在 $\mathfrak m^2$ 上为零。事实上，若 $x_1,x_2\in\mathfrak m$，则
$f'_R(x_1x_2)-x_1x_2\in\mathfrak m^{r+1}$，因为由条件 a) 有
\[
f'_R(x_1)=x_1+\alpha_1,\quad\alpha_1\in\mathfrak m^r,
\qquad
f'_R(x_2)=x_2+\alpha_2,\quad\alpha_2\in\mathfrak m^r,
\]
从而
\[
\begin{aligned}
f'_R(x_1x_2)&=f'_R(x_1)f'_R(x_2)=(x_1+\alpha_1)(x_2+\alpha_2)\\
&=x_1x_2+x_1\alpha_2+x_2\alpha_1+\alpha_1\alpha_2,
\end{aligned}
\]
而显然 $x_1\alpha_2+x_2\alpha_1+\alpha_1\alpha_2\in\mathfrak m^{r+1}$。

因此，取商后得到一个\Term{abelian-group}{Abel 群}\Term{homomorphism}{同态}
\[
u(f'_R):\mathfrak m/\mathfrak m^2\longrightarrow
 \mathfrak m^r/\mathfrak m^{r+1}.
\]
\Term{homomorphism}{同态} $u(f'_R)$ 实际上是 $k$-\Term{vector-space}{向量空间}的\Term{homomorphism}{同态}。为此，只需证明对每个 $x\in\mathfrak m$、$\alpha\in R$，均有
\[
f'_R(\alpha x)-\alpha x
\equiv\alpha(f'_R(x)-x)\pmod{\mathfrak m^{r+1}},
\]
\SourcePage{431}
即
\[
f'_R(\alpha x)-\alpha f'_R(x)\in\mathfrak m^{r+1}.
\tag{4.7.1}
\]
但我们有
\[
f'_R(\alpha x)-\alpha f'_R(x)=(f'_R(\alpha)-\alpha)f'_R(x),
\]
而
\[
f'_R(\alpha)-\alpha\in\mathfrak m^r\quad\text{（条件 a)）},\qquad
f'_R(x)\in\mathfrak m\quad(x\in\mathfrak m),
\]
故得 (4.7.1)。

\Term{homomorphism}{同态} $u(f'_R)$ 还是 $k[G]$-\Term{module}{模}\Term{homomorphism}{同态}。事实上，对 $x\in\mathfrak m$，有
\[
\begin{aligned}
u(f'_R)\bigl(g(x\bmod\mathfrak m^2)\bigr)
 &=(f'_R(gx)-gx)\bmod\mathfrak m^{r+1}\\
 &=g(f'_R(x)-x)\bmod\mathfrak m^{r+1}
 =g\,u(f'_R)(x\bmod\mathfrak m^2).
\end{aligned}
\]
最后，$u(f'_R)\ne0$。事实上，由条件 b)，存在 $x_0\in R$，使
$f'_R(x_0)-x_0\notin\mathfrak m^{r+1}$。但我们需要一个属于 $\mathfrak m$ 的这种元素。为得到它，将 $x_0$ 写为
\[
x_0=\frac{x}{x'},\qquad x,x'\in\mathfrak m,\quad x'\notin\mathfrak m^{r+1}.
\]
$x,x'$ 中至少有一个满足所需条件。
\SourcePage{432}
否则，假定同时有
\[
f'_R(x)-x\in\mathfrak m^{r+1},\qquad
f'_R(x')-x'\in\mathfrak m^{r+1}.
\]
那么，在 $A$ 的\Term{fraction-field}{分式域} $K$ 中有
\[
\frac{f'_R(x)}x\in1+\mathfrak m^r,\qquad
\frac{f'_R(x')}{x'}\in1+\mathfrak m^r.
\]
可以假定 $f'_R(x')\ne0$，否则 $x'$ 就是所求的元素。在此假设下，由于 $1+\mathfrak m^r$ 是一个\Term{multiplicative-group}{乘法群}，得到
\[
\frac{x'}{f'_R(x')}\in1+\mathfrak m^r,
\]
从而
\[
\frac{f'_R(x)}x\cdot\frac{x'}{f'_R(x')}\in1+\mathfrak m^r.
\]
但最后这一关系\Term{pushout}{推出}
\[
\frac{f'_R(x_0)}{x_0}\in1+\mathfrak m^r.
\]
事实上，
\[
\begin{aligned}
\frac{f'_R(x_0)}{x_0}-\frac{f'_R(x)}x\frac{x'}{f'_R(x')}
 &=\frac{f'_R(x/x')}{x/x'}-\frac{f'_R(x)}x\frac{x'}{f'_R(x')}\\
 &=\frac{x'f'_R(x/x')}{x}-\frac{f'_R(x)}x\frac{x'}{f'_R(x')}\\
 &=\frac1x\left(
 \frac{x'f'_R(x')f'_R(x/x')-x'f'_R(x)}{f'_R(x')}
 \right)=0.
\end{aligned}
\]
\SourcePage{433}
给 $k$-\Term{vector-space}{向量空间}
\[
\operatorname{Hom}_k(\mathfrak m/\mathfrak m^2,
 \mathfrak m^r/\mathfrak m^{r+1})
\]
赋予通常的 $G$-\Term{module}{模}结构后，其元素 $u(f'_R)$ 属于 $G$ 下的\Term{invariant-subspace}{不变子空间}
\[
\operatorname{Hom}_k(\mathfrak m/\mathfrak m^2,
 \mathfrak m^r/\mathfrak m^{r+1})^G.
\]
另一方面，经典事实是：\Term{vector-space}{向量空间} $\mathfrak m^i/\mathfrak m^{i+1}$ \Term{canonical}{典范}地等同于第 $i$ 次\Term{tensor-power}{张量幂}
$\bigotimes_k^i(\mathfrak m/\mathfrak m^2)$。因此，\Term{vector-space}{向量空间}
$\operatorname{Hom}_k(\mathfrak m/\mathfrak m^2,\mathfrak m^r/\mathfrak m^{r+1})$
通过“缩并”\Term{homomorphism}{同态}与第 $(r-1)$ 次\Term{tensor-power}{张量幂}
$\bigotimes_k^{r-1}(\mathfrak m/\mathfrak m^2)$ 等同：
\[
\gamma:\operatorname{Hom}_k(\mathfrak m/\mathfrak m^2,
 \mathfrak m^r/\mathfrak m^{r+1})
\longrightarrow\bigotimes_k^{r-1}(\mathfrak m/\mathfrak m^2).
\]
在\Term{isomorphism}{同构} $\gamma$ 下，左边的 $G$-\Term{module}{模}结构转化为右边由 $\mathfrak m/\mathfrak m^2$ 的 $G$-\Term{module}{模}结构诱导的 $G$-\Term{module}{模}结构：
\[
g(\xi_1\otimes\xi_2\otimes\cdots\otimes\xi_{r-1})
=g\xi_1\otimes\cdots\otimes g\xi_{r-1}.
\]
以 $\mu(f'_R)$ 表示 $u(f'_R)$ 在\Term{isomorphism}{同构} $\gamma$ 下的像。于是有
\[
\text{I)}\quad0\ne\mu(f'_R)\in\bigotimes_k^{r-1}(\mathfrak m/\mathfrak m^2);
\]
\SourcePage{434}
\[
\text{II)}\quad\mu(f'_R)\text{ 在 }G\text{ 下不变。}
\]
现在可以完成\Term{lemma}{引理} 4.7 的证明。首先，由于 $\mathfrak m/\mathfrak m^2$ 是一维\Term{vector-space}{向量空间}，有
\[
\operatorname{Aut}_k(\mathfrak m/\mathfrak m^2)=k^*.
\]
因此，$\operatorname{Im}(G\longrightarrow\operatorname{Aut}_k(\mathfrak m/\mathfrak m^2))$
是 $k^*$ 的一个\Term{finite}{有限}循环\Term{subgroup}{子群}，由一个本原 $c$ 次\Term{root-of-unity}{单位根}生成。以 $\varepsilon$ 表示这样的一个根。上述像\Term{subgroup}{子群}于是由下列\Term{automorphism}{自同构}生成：
\[
\chi:\mathfrak m/\mathfrak m^2\longrightarrow\mathfrak m/\mathfrak m^2,
\qquad\xi\longmapsto\varepsilon\xi.
\]
由于 $\bigotimes_k^{r-1}(\mathfrak m/\mathfrak m^2)$ 的元素 $\mu(f'_R)$ 在 $G$ 下\Term{invariant}{不变}，得到
\[
\begin{aligned}
\mu(f'_R)&=\chi\mu(f'_R)
 =\chi\sum_i\xi_{1i}\otimes\xi_{2i}\otimes\cdots\otimes\xi_{r-1,i}\\
 &=\sum_i\varepsilon\xi_{1i}\otimes\varepsilon\xi_{2i}
   \otimes\cdots\otimes\varepsilon\xi_{r-1,i}
 =\varepsilon^{r-1}\mu(f'_R),
\end{aligned}
\]
即 $(1-\varepsilon^{r-1})\mu(f'_R)=0$。由于 $\mu(f'_R)$ 非零（条件 I），有
$1-\varepsilon^{r-1}=0$，也就是说 $c\mid(r-1)$。
\end{proof}

\SourcePage{435}
\begin{proof}[命题 4.6 的证明]
已经指出，可以归约到 $g=e$、$\varphi=\operatorname{id}_G$。以 $G_1$ 表示\Term{homomorphism}{同态}
\[
G\longrightarrow\operatorname{Aut}_k(\mathfrak m/\mathfrak m^2)
\]
的核。由\Term{ramification}{分歧}理论，
\[
\operatorname{Card}(G_1)=p^t.
\]
因此，沿用\Term{lemma}{引理} 4.7 的记号，有
\[
\operatorname{Card}(G)=p^t c.
\]
另一方面，$c^{\operatorname{id}_G}(e)=\operatorname{Card}(G)$。所以，命题 4.6 在 $g=e$、$\varphi=\operatorname{id}_G$ 时的断言，以及一般情形的断言，均由\Term{lemma}{引理} 4.7 得出。
\end{proof}

\begin{SGAresult}{推论}{4.8}
对每个 $\ell\ne\operatorname{char}(k)$，有理数
\[
\frac{\nu(g_R^{-1}f'_R)-1}{c^\varphi(g)}
\]
在\Term{canonical}{典范}\Term{embedding}{嵌入} $\mathbf Q\hookrightarrow\mathbf Q_\ell$ 下的像是一个\Term{ell-adic-integer}{$\ell$ 进整数}。
\end{SGAresult}

\begin{SGAresult}{命题}{4.9}
若\Term{scheme}{概形} $X$ 的\Term{dimension}{维数}等于 $1$，则猜想 4.5 成立。
\end{SGAresult}
\SourcePage{436}
\begin{proof}
引入\Term{endomorphism}{自同态} $\psi:G\longrightarrow G$，定义为
\[
\psi=\operatorname{int}_g\circ\varphi,\qquad
\operatorname{int}_g(\xi)=g^{-1}\xi g.
\]
容易看出：

\textup{i)} $\psi(G_{x'})\subset G_{g^{-1}f'x'}$，其中 $G_{x'}$ 是点 $x'$ 的\Term{stabilizer}{稳定子群}；

\textup{ii)} $c^\varphi(g)=\operatorname{Card}(G^\psi)$，其中
$G^\psi=\{\xi\mid\psi(\xi)=\xi\}$；

\textup{iii)} $G^\psi$ \Term{action}{作用}于 $X_x'{}^{g^{-1}f'}$；

\textup{iv)} 对每个 $h\in G^\psi$，有
$\nu_{x'}(g^{-1}f')=\nu_{hx'}(g^{-1}f')$；

\textup{v)} $(g^{-1}f')(g_1x)=\psi(g_1)(g^{-1}f')(x)$。

因而可以把 $X_x'{}^{g^{-1}f'}$ \Term{resolution}{分解}为\Term{orbit}{轨道}的并，并得到
\[
\begin{aligned}
\sum_{x'\in X_x'{}^{g^{-1}f'}}\bigl(\nu_{x'}(g^{-1}f')-1\bigr)
=\sum_{x'\in X_x'{}^{g^{-1}f'}/G^\psi}
 \frac{\nu_{x'}(g^{-1}f')-1}{\operatorname{Card}(G_{x'}^\psi)}
 \operatorname{Card}(G^\psi),
\end{aligned}
\tag{4.9.1}
\]
\[
\nu_{x'}(g^{-1}f')=\nu\bigl((g^{-1}f')_{\mathscr O_{x'}}\bigr).
\]
也就是说，
\SourcePage{437}
\[
S_x(f',f,G,\varphi)(g)
=\sum_{x'\in X_x'{}^{g^{-1}f'}/G^\psi}
 \frac{\nu_{x'}(g^{-1}f')-1}{\operatorname{Card}(G_{x'}^\psi)}.
\tag{4.9.2}
\]
命题遂由命题 4.6 得出。
\end{proof}

\SGAsection{5}{Nielsen--Wecken 型公式的推广}
参见文献 [1]。
\begin{SGAresult}{命题}{5.1}
沿用第 1、2、3、4 节的假设与记号，再假定 $Y'$ 是 $X'$ 的一个\Term{finite}{有限}\Term{subset}{子集}。若 $\ell$ 是异于 $k$ 的\Term{characteristic}{特征}的素数，则：

\textup{a)} \Term{twisted-central-function}{$\varphi$ 中心函数}
\[
\sum_{x\in Y^f}S_{x,\ell}(f,f',G,\varphi)
\]
的值都是\Term{ell-adic-integer}{$\ell$ 进整数}；

\textup{b)} 下式成立：
\[
\begin{aligned}
\operatorname{Tr}_{\mathbf Z_\ell[G],\varphi}
 (f'_{\mathbb H_{U',X'}^\vee(\ell)})
={}&\sum_{x\in U^f}\nu_x(f)\operatorname{NW}_{f'}^{G,\varphi}(x)\\
&+\sum_{x\in Y^f}S_{x,\ell}(f,f',G,\varphi).
\end{aligned}
\tag{5.1.1}
\]
\end{SGAresult}
\begin{proof}
显然，只要将两边都视为取值于 $\mathbf Q_\ell$ 的函数并证明等式 b)，就已经得到 a)，因为
$\operatorname{NW}_{f'}^{G,\varphi}(x)$ 与
$\operatorname{Tr}_{\mathbf Z_\ell[G],\varphi}(f'_{\mathbb H_{U',X'}^\vee(\ell)})$
都是取值于 $\mathbf Z_\ell$ 的函数。
\SourcePage{438}
首先，对每个 $n\geq1$ 模 $\ell^n$ \Term{reduced}{约化}，并使用第 XI 讲 7.2.3.6，可知对每个 $g\in G$ 都有
\[
\begin{aligned}
c^\varphi(g)\operatorname{Tr}_{\mathbf Z_\ell[G],\varphi}
 (f'_{\mathbb H_{U',X'}^\vee(\ell)})(g)
 &=\operatorname{Tr}_{\mathbf Z_\ell}
   (g_{X'}^{-1}f')_{\mathbb H_{U',X'}^\vee(\ell)}\\
 &=\operatorname{Tr}_{\mathbf Z_\ell}
   (g_{X'}^{-1}f')_{\mathbb H_{U',X'}(\ell)}.
\end{aligned}
\tag{5.2}
\]
各行\Term{exact}{正合}的\Term{commutative-diagram}{交换图} (1.3) 给出（第 XI 讲第 4 节）
\[
\begin{aligned}
\operatorname{Tr}_{\mathbf Z_\ell}
 (g_{X'}^{-1}f')_{\mathbb H_{U',X'}(\ell)}
={}&\operatorname{Tr}_{\mathbf Z_\ell}
 (g_{X'}^{-1}f')_{\mathrm R\Gamma_{X'}(\mathbf Z_{\ell,X'})}\\
&-\operatorname{Tr}_{\mathbf Z_\ell}
 (g_{Y'}^{-1}f'_{Y'})_{\mathrm R\Gamma_{Y'}(\mathbf Z_{\ell,Y'})}.
\end{aligned}
\tag{5.3}
\]
(5.3) 右边的两项可以用经典 Lefschetz 公式求出（参见\SGAsourcegap{原 PDF 第 438 页此处“cf.”之后的引文编号为空白。}）。事实上，最后一项只涉及一个平凡的 Lefschetz 公式：
\[
\operatorname{Tr}_{\mathbf Z_\ell}
 (g_{X'}^{-1}f')_{\mathrm R\Gamma_{X'}(\mathbf Z_{\ell,X'})}
=\sum_{x\in X^f}\sum_{x'\in X_x'{}^{g_{X'}^{-1}f'}}
 \nu_{x'}(g_{X'}^{-1}f'),
\tag{5.4}
\]
\[
\operatorname{Tr}_{\mathbf Z_\ell}
 (g_{Y'}^{-1}f'_{Y'})_{\mathrm R\Gamma_{Y'}(\mathbf Z_{\ell,Y'})}
=\text{$g_{Y'}^{-1}f'_{Y'}$ 的不动点个数。}
\tag{5.5}
\]
因此
\[
\begin{aligned}
\operatorname{Tr}_{\mathbf Z_\ell}
 (g_{X'}^{-1}f')_{\mathbb H_{U',X'}(\ell)}
={}&\sum_{x\in U^f}\sum_{x'\in X_x'{}^{g_{X'}^{-1}f'}}
 \nu_{x'}(g_{X'}^{-1}f')\\
&-\sum_{x\in Y^f}\sum_{x'\in X_x'{}^{g_{X'}^{-1}f'}}
 \bigl(1-\nu_{x'}(g_{X'}^{-1}f')\bigr).
\end{aligned}
\tag{5.6}
\]
\SourcePage{439}
然而，对给定的 $x\in U^f$，\Term{set}{集合} $X_x'{}^{g_{X'}^{-1}f'}$ \Term{nonempty}{非空}当且仅当
$\operatorname{NW}_{f'}^{G,\varphi}(x)=1$。再由
\[
\operatorname{Card}X_x'{}^{g_{X'}^{-1}f'}=c^\varphi(g),\qquad x\in U,
\tag{5.7}
\]
\[
\nu_{x'}(g_{X'}^{-1}f')=\nu_x(f),\qquad x\in U,
\tag{5.8}
\]
关系式 (5.2)、(5.6) 给出
\[
\begin{aligned}
\operatorname{Tr}_{\mathbf Z_\ell[G],\varphi}
 (f'_{\mathbb H_{U',X'}^\vee(\ell)})(g)
={}&\sum_{x\in U^f}\nu_x(f)\operatorname{NW}_{f'}^{G,\varphi}(x)(g)\\
&-\sum_{x\in Y^f}\sum_{x'\in X_x'{}^{g_{X'}^{-1}f'}}
 \frac{1-\nu_{x'}(g_{X'}^{-1}f')}{c^\varphi(g)},
\end{aligned}
\tag{5.9}
\]
这正是所需证明的公式。
\end{proof}

\textbf{5.10. 特殊情形。}
\textup{a)} \Term{unramified}{非分歧}情形：$Y'=\varnothing$，即 $Y=\varnothing$。此时 (5.1.1) 成为
\[
\operatorname{Tr}_{\mathbf Z_\ell[G],\varphi}
 (f'_{\mathbb H_{U',X'}^\vee(\ell)})
=\sum_{x\in X^f}\nu_x(f)\operatorname{NW}_{f'}^{G,\varphi}(x).
\tag{5.11}
\]
\textup{b)} \Term{multiplicity}{重数}为 $1$ 的情形：
\[
\nu_{x'}(g_{X'}^{-1}f')=1
\quad\text{对所有 }x'\in Y'{}^{g_{Y'}^{-1}f'_{Y'}}\text{ 及所有 }g\in G.
\]
此时公式 (5.1.1) 归结为公式 (5.11)。特别地，若
$\dim X=1$，且对每个 $y\in Y$ 都有 $\nu_x(f)=1$，便属于这种情形。

\SourcePage{440}
\SGAsection{6}{应用于 Lefschetz 公式}
\textbf{6.1.} 沿用第 1 节的记号与条件。\footnote{$X^f$ 表示 $X$ 的闭\Term{fixed-point}{不动点}的\Term{set}{集合}。}
再设 $A$ 是一个交换的 \Term{noetherian}{Noether} $\Lambda_n$-\Term{algebra}{代数}，$F$ 是 $X$ 上的一个由\Term{perfect-module}{完美} $A$-\Term{module}{模}组成的\Term{constructible-sheaf}{可构造层}（相对于\Term{etale-topology}{平展拓扑}），并设
\[
u:f^*(F)\longrightarrow F
\]
为 $A$-\Term{module}{模}\Term{sheaf}{层}的一个\Term{homomorphism}{同态}。

特别地，得到二元组 $(X,F)$ 的一个\Term{endomorphism}{自同态}
\[
(u,f):(X,F)\longrightarrow(X,F),
\]
从而得到 $A$-\Term{module}{模}\Term{complex}{复形}的一个\Term{endomorphism}{自同态}
\[
\mathrm R\Gamma_X(u,f):\mathrm R\Gamma_X(F)\longrightarrow\mathrm R\Gamma_X(F).
\]
依假设，\Term{scheme}{概形} $X$ 在 $k$ 上\Term{proper}{固有}，而\Term{sheaf}{层} $F$ 是\Term{constructible-sheaf}{可构造}且\Term{torsion-sheaf}{挠}的，因此\Term{finiteness-theorem}{有限性定理}（SGA 4，第 XVII 讲 5）保证 $\mathrm R\Gamma_X(F)$ 是 $A$-\Term{module}{模}的\Term{perfect-complex}{完美复形}。所以\Term{trace}{迹}
$\operatorname{Tr}_A(\mathrm R\Gamma_X(u,f))$ 有定义，且
\[
\operatorname{Tr}_A(\mathrm R\Gamma_X(u,f))\in A.
\]
我们的目标是在一些附加假设下，用 $f$ 的\Term{fixed-point}{不动点}所关联的\Term{local-invariant}{局部不变量}给出这个\Term{trace}{迹}的显式公式。

\SourcePage{441}
\textbf{6.2. \Term{local-term}{局部项}的定义。}
假定 $X'$ 是一维\Term{irreducible}{不可约}\Term{scheme}{概形}，且\Term{sheaf}{层} $F'=p^*(F)$ 在 $U'$ 上的\Term{restriction}{限制}为\Term{constant-sheaf}{常值层}。

若 $f'$ 的\Term{fixed-point}{不动点}个数\Term{finite}{有限}——下文始终作此假定——则总可以假设 $X^f\subset Y^f$。

\Term{sheaf}{层} $F'|_{U'}$ 因而由一个带有 $G$ \Term{action}{作用}的 $A$-\Term{module}{模} $M$ 确定：
\[
F'|_{U'}\simeq M_{U'}.
\]
事实上，以 $\eta$（相应地以 $\eta'$）表示 $X$（相应地 $X'$）的\Term{generic-point}{泛点}，以 $\bar\eta$（相应地 $\bar\eta'$）表示 $\eta$（相应地 $\eta'$）上的一个\Term{geometric-point}{几何点}，则有
\[
p(\eta')=\eta,\qquad F_{\bar\eta}\simeq F'_{\bar\eta'}.
\]
\Term{complex}{复形} $F'$ 与 $M$ \Term{isomorphism}{同构}，所以 $M$ 是一个\Term{perfect-module}{完美} $A$-\Term{module}{模}。

另一方面，若\Term{morphism}{态射} $f'$ 不是常值的，则它固定 $X'$ 的\Term{generic-point}{泛点}。因此，二元组 $(X',F')$ 的\Term{morphism}{态射} $(p^*(u),f')$ 在\Term{stalk}{茎} $F'_{\bar\eta'}$ 上诱导一个\Term{morphism}{态射}
\[
u'_{\bar\eta'}:F'_{\bar\eta'}\longrightarrow F'_{\bar\eta'},
\]
从而显然得到一个 $A$-\Term{module}{模}\Term{morphism}{态射}
\[
u_M:M\longrightarrow M.
\]
\SourcePage{442}
容易看出，下列关系成立：
\[
u_M(\varphi(g)m)=gu_M(m)\qquad(\forall m\in M,\ \forall g\in G).
\tag{6.2.1}
\]
因此，$M$ 相对于 $A$ 的\Term{dual}{对偶} $M^\vee$ 带有 $A[G]$-\Term{module}{模}结构，并有一个\Term{endomorphism}{自同态}
\[
u_M^\vee:M^\vee\longrightarrow M^\vee,
\]
满足
\[
u_M^\vee(gx)=\varphi(g)u_M^\vee(x)
\qquad(\forall g\in G,\ \forall x\in M^\vee).
\]
此外，$M^\vee$ 是一个\Term{perfect-module}{完美} $A$-\Term{module}{模}，所以对每个 $g\in G$，\Term{trace}{迹} $\operatorname{Tr}_A(g^{-1}u_M^\vee)$ 都有定义。

设 $x\in X^f$ 是 $f$ 的一个闭\Term{fixed-point}{不动点}。\Term{morphism}{态射} $u$ 诱导 $F$ 在 $x$ 处的\Term{stalk}{茎} $F_x$ 的一个\Term{endomorphism}{自同态}（我们将 $X$ 的\Term{closed-point}{闭点}与 $X(k)$ 的点等同）：
\[
u_x:F_x\longrightarrow F_x.
\]
\Term{trace}{迹} $\operatorname{Tr}_A(u_x)$ 给出了与\Term{fixed-point}{不动点} $x$ 相关的第一个\Term{local-invariant}{局部不变量}。

再由下式引入另一个\Term{local-invariant}{局部不变量} $\alpha_x(f,u,f',G,\varphi)\in A$：
\[
\alpha_x(f,u,f',G,\varphi)
=\sum_{g\in G_{\natural,\varphi}}
 S_{x,\ell^n}(f,f',G,\varphi)(g)\operatorname{Tr}_A(g^{-1}u_M^\vee).
\tag{6.2.2}
\]
\SourcePage{443}
这里函数
\[
S_{x,\ell^n}(f,f',G,\varphi):G\longrightarrow\Lambda_n
\]
由函数 $i_\ell\circ S_x(f,f',G,\varphi)$（参见第 4 节）模 $\ell^n$ \Term{reduced}{约化}得到，即
\[
S_{x,\ell^n}(f,f',G,\varphi)(g)
=(i_\ell\circ S_x(f,f',G,\varphi))(g)\bmod\ell^n\mathbf Z_\ell.
\]

\textbf{6.3. 一个 Lefschetz 公式。}
在 5.1、6.1、6.2 的假设与记号下，有公式
\[
\operatorname{Tr}_A(\mathrm R\Gamma_X(u,f))
=\sum_{x\in X^f}\tau_x(f,u,f',G,\varphi),
\tag{6.3.1}
\]
其中定义
\[
\tau_x(f,u,f',G,\varphi)
=\operatorname{Tr}_A(u_x)-\alpha_x(f,u,f',G,\varphi).
\tag{6.3.2}
\]
\begin{proof}
由于\Term{subscheme}{子概形} $Y$ 在 $f$ 下\Term{stable}{稳定}，\Term{morphism}{态射}
\[
(u,f):(X,F)\longrightarrow(X,F)
\]
诱导一个\Term{distinguished-triangle}{可区别三角形}的\Term{endomorphism}{自同态}：
\[
\begin{tikzcd}[column sep=large,row sep=large]
&\mathrm R\Gamma_Y(F|_Y)
 \arrow[dl,"{+1}"']
 \arrow[loop above,"{\mathrm R\Gamma_Y(u|_Y,f|_Y)}"]&\\
\mathrm R\Gamma_X(F_{U,X})
 \arrow[rr]
 \arrow[loop left,"{\mathrm R\Gamma_X((u,f)_{U,X})}"']
&&\mathrm R\Gamma_X(F)
 \arrow[ul]
 \arrow[loop right,"{\mathrm R\Gamma_X(u,f)}"]
\end{tikzcd}
\tag{6.3.3}
\]
于是有
\SourcePage{444}
\[
\begin{aligned}
\operatorname{Tr}_A(\mathrm R\Gamma_X(u,f))
={}&\operatorname{Tr}_A(\mathrm R\Gamma_X((u,f)_{U,X}))\\
&+\operatorname{Tr}_A(\mathrm R\Gamma_Y(u|_Y,f|_Y)).
\end{aligned}
\tag{6.3.4}
\]
由初等 Lefschetz 公式，并利用总可假设 $X^f=Y^f$ 这一事实，(6.3.4) 的最后一项可明确写为
\[
\operatorname{Tr}_A(\mathrm R\Gamma_Y(u|_Y,f|_Y))
=\sum_{x\in X^f}\operatorname{Tr}_A(u_x).
\tag{6.3.5}
\]
另一方面，有
\[
\mathrm R\Gamma_X((u,f)_{U,X})
\simeq\mathrm R\Gamma^G\mathrm R\Gamma_{X'}((u',f')_{U',X'}),
\qquad u'=p^*(u),
\tag{6.3.6}
\]
\[
\begin{gathered}
F_{U',X'}\simeq\Lambda_{n,U',X'}\otimes_{\Lambda_n}M,\\
(u',f')_{U',X'}=f'_{\Lambda_{n,U',X'}}\otimes_{\Lambda_n}M.
\end{gathered}
\tag{6.3.7}
\]
因此，由 \Term{kunneth-formula}{Künneth 公式}，
\[
\mathrm R\Gamma_{X'}((u',f')_{U',X'})
\simeq\mathrm R\Gamma_{X'}(f'_{\Lambda_{n,U',X'}})
 \mathbin{\overset{\mathrm L}{\otimes}}_{\Lambda_n}u_M.
\tag{6.3.8}
\]
由 (6.3.8) \Term{pushout}{推出}
\[
\begin{aligned}
\mathrm R\Gamma_{X'}((u',f')_{U',X'})^\vee
\simeq\bigl(\mathrm R\Gamma_{X'}(f'_{\Lambda_{n,U',X'}})^\vee
 \mathbin{\overset{\mathrm L}{\otimes}}_{\Lambda_n}A\bigr)
 \mathbin{\overset{\mathrm L}{\otimes}}_A u_M^\vee.
\end{aligned}
\tag{6.3.9}
\]
根据 (6.3.6)，得到
\[
\mathrm R\Gamma_X((u,f)_{U,X})
=\mathrm R\Gamma_{X'}(f'_{\Lambda_{n,U',X'}})^\vee
 \mathbin{\overset{\mathrm L}{\otimes}}_{A[G]}u_M^\vee.
\tag{6.3.10}
\]
再由一般公式（参见第 XI 讲 7.5.9），
\[
\operatorname{Tr}_A((u,f)_{U,X})
=\sum_{g\in G_{\natural,\varphi}}
 \operatorname{Tr}_{\Lambda_n[G],\varphi}
 (f'_{\mathbb H_{n,U',X'}^\vee})(g)
 \operatorname{Tr}_A(g^{-1}u_M^\vee).
\tag{6.3.11}
\]
而\Term{trace}{迹}
$\operatorname{Tr}_{\Lambda_n[G],\varphi}(f'_{\mathbb H_{n,U',X'}^\vee})$
可用 Nielsen--Wecken 公式（参见命题 5.1 b)）计算；当然须模 $\ell^n\mathbf Z_\ell$ \Term{reduced}{约化}。
\SourcePage{445}
考虑关系式 $X^f=Y^f$，得到
\[
\operatorname{Tr}_{\Lambda_n[G],\varphi}
 (f'_{\mathbb H_{n,U',X'}^\vee})(g)
=\sum_{x\in X^f}S_{x,\ell^n}(f,f',G,\varphi)(g).
\tag{6.3.12}
\]
将 (6.3.12)、(6.3.11)、(6.3.5)、(6.3.4) 合并，即得所述公式。
\end{proof}

\textbf{6.4. 关于\Term{local-term}{局部项}的注记。}
\textup{a)} 容易看出，若 $g\in G$，则有
\[
\alpha_x(f,u,g^{-1}f',G,\varphi_g)
=\alpha_x(f,u,f',G,\varphi),
\tag{6.4.1}
\]
其中定义
\[
\varphi_g(\xi)=(\operatorname{int}_g\circ\varphi)(\xi)
=g^{-1}\varphi(\xi)g\qquad(\forall\xi\in G).
\tag{6.4.2}
\]
由 (6.4.1)，将 $f'$ 换为 $g^{-1}f'$，将 $\varphi$ 换为 $\varphi_g$ 后，总可以归约到至少存在一个点 $x'\in X_x'{}^{f'}$ 的情形。

\textup{b)} 取 $x'\in X_x'{}^{f'}$，以 $G_{x'}$ 表示 $x'$ 的\Term{stabilizer}{稳定子群}。显然，$G_{x'}$ 在 $\varphi$ 下\Term{stable}{稳定}，从而得到一个\Term{endomorphism}{自同态}
\[
\varphi_{x'}=\varphi|_{G_{x'}}:G_{x'}\longrightarrow G_{x'}.
\]
以 $X'/G_{x'}$ 表示按 $G_{x'}$ 确定的\Term{equivalence-relation}{等价关系}对 $X'$ 取商得到的\Term{quotient-scheme}{商概形}；以 $F'/G_{x'}$ 表示由 $X'$ 的\Term{endomorphism}{自同态} $f'$ 按同一\Term{equivalence-relation}{等价关系}得到的 $X'/G_{x'}$ 的\Term{endomorphism}{自同态}；以 $\bar x'$ 表示 $x'$ 模此\Term{equivalence-relation}{等价关系}的类。
\SourcePage{446}
\Term{canonical-projection}{典范投影}
\[
p:X'\longrightarrow X'/G=X
\]
可如下\Term{resolution}{分解}：
\[
p:X'\xrightarrow{p_1}X'/G_{x'}\xrightarrow{p_2}X=X'/G.
\]
对\Term{homomorphism}{同态} $u$ 沿 $p_2$ 取\Term{inverse-image}{逆像}，得到一个\Term{homomorphism}{同态}
\[
v=p_2(u):(f'/G_{x'})^*(p_2^*(F))\longrightarrow p_2^*(F).
\]
容易看出
\[
\alpha_{\bar x'}(f'/G_{x'},u,f',G_{x'},\varphi_{x'})
=\alpha_x(f,u,f',G,\varphi)\qquad(x=p(x')).
\tag{6.4.3}
\]
因而总可以归约到“\Term{totally-ramified}{全分歧}”情形，也就是点 $x$ 的\Term{fiber}{纤维}仅含一个点的情形。

\textup{c)} “\Term{tamely-ramified}{温和分歧}”情形。假定存在 $x'\in p^{-1}(x)$，使
\[
f'(x')=x',\qquad p(x')=x,
\]
并且群 $G_{x'}$ 的阶与域 $k$ 的\Term{characteristic}{特征}互素；总可以归约到存在这样的 $x'$ 的情形（参见 a)）。此时有
\[
\alpha_x(f,u,f',G,\varphi)
=S_{x,\ell^n}(f,f',G,\varphi)(e)\operatorname{Tr}_A(u_M^\vee).
\tag{6.4.4}
\]
由 b)，可以假定 $G=G_{x'}$。于是公式 (6.4.4) 由下述\Term{lemma}{引理}得出。

\SourcePage{447}
\textit{〔原 PDF 此页重复上一页的正文；以下依原样保留。〕}

\Term{canonical-projection}{典范投影}
\[
p:X'\longrightarrow X'/G=X
\]
可如下\Term{resolution}{分解}：
\[
p:X'\xrightarrow{p_1}X'/G_{x'}\xrightarrow{p_2}X=X'/G.
\]
对\Term{homomorphism}{同态} $u$ 沿 $p_2$ 取\Term{inverse-image}{逆像}，得到一个\Term{homomorphism}{同态}
\[
v=p_2(u):(f'/G_{x'})^*(p_2^*(F))\longrightarrow p_2^*(F).
\]
容易看出
\[
\alpha_{\bar x'}(f'/G_{x'},u,f',G_{x'},\varphi_{x'})
=\alpha_x(f,u,f',G,\varphi)\qquad(x=p(x')).
\tag{6.4.3}
\]
因而总可以归约到“\Term{totally-ramified}{全分歧}”情形，也就是点 $x$ 的\Term{fiber}{纤维}仅含一个点的情形。

\textup{c)} “\Term{tamely-ramified}{温和分歧}”情形。假定存在 $x'\in p^{-1}(x)$，使
\[
f'(x')=x',\qquad p(x')=x,
\]
并且群 $G_{x'}$ 的阶与域 $k$ 的\Term{characteristic}{特征}互素；总可以归约到存在这样的 $x'$ 的情形（参见 a)）。此时有
\[
\alpha_x(f,u,f',G,\varphi)
=S_{x,\ell^n}(f,f',G,\varphi)(e)\operatorname{Tr}_A(u_M^\vee).
\tag{6.4.4}
\]
由 b)，可以假定 $G=G_{x'}$。于是公式 (6.4.4) 由下述\Term{lemma}{引理}得出。

\SourcePage{448}
\begin{SGAresult}{\Term{lemma}{引理}}{6.4.5}
若下列关系成立：
\[
\nu_{x'}(f')>1,\qquad\nu_{x'}(g^{-1}f')>1,
\]
则 $g$ 属于第一\Term{ramification-group}{分歧群} $G_1$（参见 4.7）：
\[
g\in G_1,
\]
也就是说，$g$ 在\Term{vector-space}{向量空间} $\mathfrak m_{x'}/\mathfrak m_{x'}^2$ 上诱导\Term{identity-map}{恒等映射}。
\end{SGAresult}
\begin{proof}
设 $\widehat{\mathscr O}_{x'}=k[[T]]$。由 $g^{-1}$ 诱导的 $\mathfrak m/\mathfrak m^2$ 的\Term{automorphism}{自同构}，由满足下列条件的\Term{root-of-unity}{单位根} $\varepsilon(g^{-1})\in k$ 确定：
\[
g^{-1}T=\varepsilon(g^{-1})T\pmod{\mathfrak m^2},
\]
其中 $\mathfrak m$ 是环 $k[[T]]$ 的\Term{maximal-ideal}{极大理想}。

另一方面，由假设 $\nu_{x'}(f')>1$ 得到
\[
f'(T)=T+a_\nu T^\nu+a_{\nu+1}T^{\nu+1}+\cdots,\qquad\nu\geq2.
\]
因此
\[
g^{-1}f'(T)=\varepsilon(g^{-1})T+b_\nu T^\nu+b_{\nu+1}T^{\nu+1}+\cdots,
\]
而假设 $\nu(g^{-1}f')>1$ \Term{pushout}{推出} $\varepsilon(g^{-1})=1$。证毕。
\end{proof}

\SourcePage{449}
\textup{d)} \Term{sheaf}{层} $F$ 在\Term{endomorphism}{自同态} $f$ 的\Term{fixed-point}{不动点}处不\Term{ramification}{分歧}的情形。

此时，\Term{local-term}{局部项} $\tau_x(f,u,f',G,\varphi)$ 由下式给出：
\[
\tau_x(f,u,f',G,\varphi)=\nu_x(f)\operatorname{Tr}_A(u_x),
\tag{6.4.6}
\]
因而它只依赖于二元组 $(f,u)$。

事实上，此时有
\[
\begin{aligned}
{}&S_{x,\ell^n}(f,f',G,\varphi)(g)\\
&\quad=\begin{cases}
0,&\begin{aligned}
 {}&\text{存在 }x'\in X'\text{ 使 }(g^{-1}f')(x')=x',\\[-2pt]
 &\text{等价地，}g\text{ 的 }\varphi\text{ 共轭类等于}\\[-2pt]
 &\operatorname{NW}_{f'}^{G,\varphi}(x)\text{（参见第 3 节）；}
\end{aligned}\\
\nu_x(f)-1,&\text{否则。}
\end{cases}
\end{aligned}
\tag{6.4.7}
\]
特别地，由此得到
\[
S_{x,\ell^n}(f,f',G,\varphi)(g)=0
\quad\text{若 }g\notin\operatorname{NW}_{f'}^{G,\varphi}(x).
\tag{6.4.8}
\]
于是得出
\[
\begin{aligned}
\tau'_x(f,u,f',G,\varphi)
 &=\nu_x(f)\operatorname{Tr}_A(g^{-1}u_M^\vee)
  =\nu_x(f)\operatorname{Tr}_A(g^{-1}u')_{x'}\\
 &=\nu_x(f)\operatorname{Tr}_A(u_x).
\end{aligned}
\tag{6.4.9}
\]
\textit{〔原文 (6.4.7) 的两种取值次序与紧接其后的 (6.4.8) 不一致；此处保留原式。〕}

\SourcePage{450}
\SGAsection{7}{关于 Lefschetz 公式成立条件的说明}
\textbf{7.1.} 当然，在实际应用中，出发点是 $k$ 上的一个\Term{proper}{固有}\Term{scheme}{概形} $X$ 及一个\Term{endomorphism}{自同态}
\[
(u,f):(X,F)\longrightarrow(X,F),
\]
满足：
\begin{itemize}
\item $F$ 是 $A$-\Term{module}{模}的一个\Term{constructible-sheaf}{可构造层}；
\item $f$ 的\Term{fixed-point}{不动点}\Term{set}{集合} $X^f$ \Term{finite}{有限}。
\end{itemize}
问题是在这些条件下定义\Term{local-term}{局部项} $\tau_x(u,f)$，$x\in X^f$，使得
\[
\operatorname{Tr}_A(\mathrm R\Gamma_X(u,f))
=\sum_{x\in X^f}\tau_x(u,f).
\tag{*}
\]
利用整体与\Term{local-scheme}{局部}两个\Term{duality-theorem}{对偶定理}，Verdier 已经定义了使公式 $(*)$ 成立的\Term{local-term}{局部项} $\tau_x(u,f)$（参见第 III 讲）。遗憾的是，似乎很难用已知的\Term{local-invariant}{局部不变量}显式算出 Verdier 所定义的 $\tau_x(u,f)$；这正是前一节提出那些定义的意义所在。\footnote{不过，关于 (6.3.1) 的\Term{local-term}{局部项}与 Verdier 的\Term{local-term}{局部项}的比较，参见第 III B 讲。}

另一方面，后面这些定义需要对 $F$、$f$、$u$ 作附加假设，即要求存在满足 5.1、6.1、6.2 所需性质的 $U,U',G,f',\varphi$。

\SourcePage{451}
\begin{SGAresult}{命题}{7.2}
设 $X$ 是一维\Term{smooth}{光滑}\Term{scheme}{概形}，$F$ \Term{constructible-sheaf}{可构造}，并设
\[
(u,f):(X,F)\longrightarrow(X,F)
\]
为一个\Term{endomorphism}{自同态}，使得\Term{sheaf}{层}\Term{morphism}{态射}
\[
u:f^*(F)\longrightarrow F
\]
是\Term{isomorphism}{同构}。在这些条件下，存在满足 5.1、6.1、6.2 的条件的 $U,U',G,f',\varphi$。
\end{SGAresult}
\begin{proof}
设 $U$ 是 $X$ 的使 $F$ 在其上\Term{locally-constant-sheaf}{局部常值}的最大\Term{open-subscheme}{开子概形}，以 $Y$ 表示其补集。\Term{set}{集合} $Y$ 在 $f$ 下\Term{stable}{稳定}：
\[
f(Y)\subset Y.
\]
事实上，反设存在 $y\in Y$，使 $f(y)\in U$。那么\Term{sheaf}{层} $f^*(F)$ 在 $y$ 的某个\Term{neighborhood}{邻域}内\Term{locally-constant-sheaf}{局部常值}；利用\Term{isomorphism}{同构} $u$，得到 $F$ 在 $y$ 的某个\Term{neighborhood}{邻域}内也\Term{locally-constant-sheaf}{局部常值}，矛盾。

若 $f$ 为常值，命题的证明是平凡的。因此假定 $f$ 不是常值。若 $\eta$ 是 $X$ 的\Term{generic-point}{泛点}，则有 $f(\eta)=\eta$。
\SourcePage{452}
设 $i_\eta:\operatorname{Spec}(k(\eta))\hookrightarrow X$ 为\Term{canonical}{典范}\Term{embedding}{嵌入}，$F_\eta=i_\eta^*(F)$，并以
\[
u_\eta:f_\eta^*(F_\eta)\longrightarrow F_\eta
\]
表示由 $u$ 诱导的\Term{sheaf}{层}\Term{morphism}{态射}（实际上是\Term{isomorphism}{同构}）。这里 $f_\eta$ 当然是 $f$ 在 $\operatorname{Spec}(k(\eta))$ 上的\Term{restriction}{限制}。

将一切改用 $H$-\Term{module}{模}来表述，其中
\[
H=\operatorname{Gal}(k(\eta)_s/k(\eta)),\qquad
k(\eta)_s\text{ 是 }k(\eta)\text{ 的一个可分闭包},
\]
得到\Term{isomorphism}{同构}
\[
H\xrightarrow{\psi}H,\qquad
F_{\bar\eta}\xrightarrow{\underline u_\eta}F_{\bar\eta},
\]
其中 $F_{\bar\eta}$ 带有 $H$-\Term{module}{模}结构，且满足
\[
\underline u_\eta(gx)=\psi(g)\underline u_\eta x,
\qquad g\in H,\quad x\in F_{\bar\eta}.
\]
若 $H_1$ 是 $H$ 中使 $F_{\bar\eta}^{H_1}=F_{\bar\eta}$ 的最大的\Term{finite-index}{有限指数}\Term{normal-subgroup}{正规子群}，则 $\psi(H_1)=H_1$。这样的\Term{subgroup}{子群}存在，因为\Term{sheaf}{层} $F_\eta$ 是\Term{constructible-sheaf}{可构造}的。

设 $K\supset k(\eta)$ 是与\Term{subgroup}{子群} $H_1$ 对应的 $k(\eta)$ 的\Term{galois-extension}{Galois 扩张}，$X'$ 是 $X$ 在 $K$ 中的\Term{normalization}{正规化}。显然，可以定义\Term{endomorphism}{自同态}
\[
f':X'\longrightarrow X'
\]
以及\Term{endomorphism}{自同态}
\[
H/H_1=G\xrightarrow{\varphi}G,
\]
使得系统 $U,U',f',G,\varphi$ 满足所需条件。
\end{proof}

\SourcePage{453}
\begin{SGAresult}{例}{7.3}
设 $X$ 是\Term{finite-field}{有限域} $\mathbf F_p$ 上的一个\Term{prescheme}{预概形}，$p$ 是素数，$\operatorname{fr}_X:X\longrightarrow X$ 是 Frobenius \Term{endomorphism}{自同态}（参见第 XIV 讲），$F$ 是 $X$ 上的一个 $\Lambda$-\Term{module}{模}\Term{sheaf}{层}。

若 $Y,S$ 是 $\mathbf F_p$ 上的\Term{scheme}{概形}，且 $Y$ 是一个 $S$-\Term{scheme}{概形}，则以
\[
\operatorname{Fr}_{Y/S}:Y\longrightarrow Y^{(p/S)}
\]
表示相对 Frobenius \Term{morphism}{态射}（参见第 XIV 讲 2）。

可以得到 $X_{\text{ét}}$ 上的\Term{sheaf}{层}\Term{morphism}{态射}
\[
(\operatorname{fr}_X)_*(F)\longrightarrow F,
\]
其定义为：对每个\Term{etale}{平展} $X$-\Term{prescheme}{预概形} $U$，令
\[
\begin{aligned}
(\operatorname{fr}_X)_*(F)(U)
 &=F(\operatorname{fr}_X^{-1}(U))=F(U^{(p/X)})\\
 &\xrightarrow{F(\operatorname{Fr}_{U/X})}F(U).
\end{aligned}
\]
此\Term{morphism}{态射}实际上是一个\Term{isomorphism}{同构}。其逆
\[
F(\operatorname{Fr}_{/X})^{-1}:F\longrightarrow(\operatorname{fr}_X)_*(F)
\]
由\Term{adjunction}{伴随}定义出一个\Term{morphism}{态射}
\[
\operatorname{Fr}_{F/X}^{*}:(\operatorname{fr}_X)^*(F)\longrightarrow F,
\]
它是一个\Term{isomorphism}{同构}。

特别地，若 $X$ 一维\Term{smooth}{光滑}，$F$ \Term{constructible-sheaf}{可构造}，则二元组
\[
(\operatorname{Fr}_{F/X}^{*},\operatorname{fr}_X):(X,F)\longrightarrow(X,F)
\]
满足命题 7.2 的条件。
\end{SGAresult}

\medskip
\noindent [1]\quad F. Wecken，\textit{Fixpunktklassen I}（\Term{fixed-point}{不动点}类 I），\textit{Math. Ann.} \textbf{117}（1941），659--671 页；
\textit{Fixpunktklassen II, III}（\Term{fixed-point}{不动点}类 II、III），\textit{Math. Ann.} \textbf{118}（1942--1943），216--234 页及 544--577 页。
